\section{A finite reduction for arbitrary mixed integer inner orders}
\label{mixed:section}

The directional-zeta report reduces two integer inner orders and the
arbitrary-depth sector in which every inner order is nonpositive. Its
Question 11 asks for a finite nested-polylogarithm reduction at arbitrary
depth with both signs present, while the outer order stays complex
\cite{RepoDirectional}. We solve that formulation. The construction uses
root filtering, shuffle products, rational partial fractions, and finite
geometric power sums. Each stage is explicit and accommodates repeated
colors. The analytic continuation and the directions in which one may
differentiate are part of the statement.

The shuffle algebra and rational manipulation of hyperlogarithms are
classical and have powerful general implementations \cite{Panzer}.
Here the desired output is a finite expression with one free complex
order, positive remaining indices, a controlled depth, and the fixed
cumulative root alphabet. The construction and its proof are given in
full, so their correctness does not depend on an external symbolic
integration package.

\subsection{Statement, alphabet, and convergence}

For integers $a_j$, positive integers $p_j$, and nonzero colors
$|X_j|\le1$, put
\begin{equation}\label{mixed:Z}
 Z(C)=\sum_{m_1,\ldots,m_d\ge1}
       \frac{\prod_{j=1}^d X_j^{m_j}m_j^{-a_j}}
            {(p_1m_1+\cdots+p_dm_d)^C}.
\end{equation}
If all colors have modulus less than one, this series is entire in
$C$. For unit colors it initially converges absolutely when, for example,
\[
              \Re C>d+\sum_j\max(0,-a_j).
\]
This sufficient half-plane follows by grouping according to the
denominator and using at most $O(n^{d-1})$ tuples of total weight $n$.
Its size is not the research claim.

We use the nested convention
\begin{equation}\label{mixed:Li}
 \Li_{s_1,\ldots,s_r}(z_1,\ldots,z_r)
  =\sum_{n_1>\cdots>n_r\ge1}
            \frac{z_1^{n_1}\cdots z_r^{n_r}}
                 {n_1^{s_1}\cdots n_r^{s_r}}.
\end{equation}
The \emph{cumulative colors} are $\xi_j=z_1\cdots z_j$.
The relative colors $z_j$ need not have modulus at most one.
If $\max_j|\xi_j|<1$, the series is locally normally convergent in
all its orders: writing $n_j-n_{j+1}$ as positive increments, with
$n_{r+1}=0$, bounds the color product by $\rho^{n_1}$ for some
$\rho<1$. Polynomial growth from the denominators is then harmless.

\begin{theorem}[Finite mixed-order normal form]\label{mixed:main}
Let $P=\#\{j:a_j>0\}$ and
\[
                \mathcal A=\bigcup_{j=1}^d\{\alpha:\alpha^{p_j}=X_j\}.
\]
The function \eqref{mixed:Z} is a finite linear combination of terms
\begin{equation}\label{mixed:output}
          \Li_{C+v,k_2,\ldots,k_r}(z_1,\ldots,z_r),
 \qquad v\in\Z,\quad k_2,\ldots,k_r\in\Z_{>0},
\end{equation}
whose cumulative colors all belong to $\mathcal A$. The coefficients
are obtained by finite rational operations in the colors and roots,
rational Bernoulli numbers, and integral shuffle multiplicities.
One may choose
\begin{equation}\label{mixed:depth}
 r\le\begin{cases}
 d,&P=d,\\
 P+1,&0<P<d,\\
 1,&P=0.
 \end{cases}
\end{equation}
If all $X_j$ are roots of unity of orders dividing $q_j$, every
cumulative color has order dividing $\lcm_j(p_jq_j)$.

The identity initially holds in the region of absolute convergence
and holds meromorphically for every $C$. If every $X_j\ne1$, both
$Z$ and every final term are entire in $C$. Every fixed derivative
in $C$ is obtained by taking the corresponding derivative with respect
to the first index
of the multiple polylogarithms in \eqref{mixed:output}.
The fixed integral orders $a_j$ are not free differentiation variables
in this formula.
\end{theorem}

The depth in \eqref{mixed:depth} is a bound for this construction;
it is not a minimal-depth or period-independence assertion. In particular,
one nonpositive factor can substantially shorten a high-depth input
when only a few of the remaining indices are positive.

\subsection{Step 1: remove weights and shuffle the positive factors}

Let the coefficient transform be
\[
 \Mcal_C\left[\sum_{n\ge0}f_nt^n\right]
             =\sum_{n\ge1}\frac{f_n}{n^C}.
\]
The constant term is omitted. The generator of \eqref{mixed:Z} is
\begin{equation}\label{mixed:generator}
 K(t)=\prod_j\Li_{a_j}(X_jt^{p_j}),\qquad Z(C)=\Mcal_C K.
\end{equation}
Root filtering gives, for every integer $a$,
\begin{equation}\label{mixed:roots}
            \Li_a(Xt^p)=p^{a-1}\sum_{\alpha^p=X}\Li_a(\alpha t).
\end{equation}
Indeed the sum of $n$th powers of the roots is zero unless $p$ divides
$n$, and is then $pX^{n/p}$. This also checks the factor $p^{a-1}$.
Expanding the product reduces the proof to unit weights and colors
in $\mathcal A$.

For a nonzero letter $\alpha$, use the one-forms
$\omega_0=dt/t$ and $\omega_\alpha=\alpha\,dt/(1-\alpha t)$,
and define iterated integrals recursively by integrating the leftmost
letter against the remaining word. The word
$0^{a-1}\alpha$ represents $\Li_a(\alpha t)$ for $a>0$.
Partitioning the product of the integration simplices by the order
of their coordinates proves the shuffle product, with multiplicities.
Thus the $P$ positive factors give finitely many words with exactly
$P$ nonzero letters. A word with zero-block indices $k_1,\ldots,k_P$
and successive nonzero letters $\beta_1,\ldots,\beta_P$ represents
\begin{equation}\label{mixed:word}
 \Li_{k_1,\ldots,k_P}
       (\beta_1t,\beta_2/\beta_1,\ldots,\beta_P/\beta_{P-1}).
\end{equation}
All the $k_i$ are positive. This identification follows either by
expanding the forms geometrically in the recursive integrals, or by
differentiating and fixing the value at $t=0$. Its cumulative colors
are exactly the word letters, all in $\mathcal A$.

The nonpositive factors form a rational function $R(t)$, since
\[
 \Li_{-M}(z)=(z\,d/dz)^M\frac{z}{1-z}\qquad(M\ge0).
\]
Each factor is bounded at infinity, so there is no positive-degree
polynomial part in the partial-fraction decomposition
\begin{equation}\label{mixed:R}
 R(t)=c_\infty+\sum_{\alpha\in\mathcal A}\sum_{h=1}^{M_\alpha}
                         \frac{c_{\alpha,h}}{(1-\alpha t)^h}.
\end{equation}
An upper bound for $M_\alpha$ is the sum of $1-a_j$ over the
nonpositive factors with color $\alpha$. Even if a numerator cancels
some poles, the valid coefficient formula remains
\begin{equation}\label{mixed:partialcoeff}
 c_{\alpha,h}=[u^{M_\alpha-h}]
                 u^{M_\alpha}R((1-u)/\alpha).
\end{equation}
This finite Taylor extraction handles repeated colors directly.

\subsection{Step 2: a rational prefactor with one free spectral order}

For a nested polylogarithm $H$ as in \eqref{mixed:word}, write its
relative colors as $\beta_1,\ldots,\beta_r$ for this subsection.
Clearly
\[
 \Mcal_C\Li_{k_1,\ldots,k_r}(\beta_1t,\beta_2,\ldots,\beta_r)
   =\Li_{C+k_1,k_2,\ldots,k_r}(\beta_1,\ldots,\beta_r).
\]
Define the rational numbers $b_{h,k}$ by
\begin{equation}\label{mixed:b}
 \binom{v+h-1}{h-1}=
       \prod_{\nu=1}^{h-1}(1+v/\nu)
       =\sum_{k=0}^{h-1}b_{h,k}v^k.
\end{equation}
Their connection to generalized harmonic numbers is explicit:
\[
 \prod_{\nu=1}^{h-1}(1+v/\nu)
   =\exp\left(\sum_{j\ge1}
          \frac{(-1)^{j+1}H_{h-1}^{(j)}}{j}v^j\right).
\]

\begin{lemma}[Rational-prefactor transform]\label{mixed:prefactor}
For $h\ge1$ and $H=\Li_{k_1,\ldots,k_r}
(\beta_1t,\beta_2,\ldots,\beta_r)$,
\begin{align}\label{mixed:prefactorformula}
 \Mcal_C\bigl[(1-\alpha t)^{-h}H\bigr]
 ={}&\Li_{C+k_1,k_2,\ldots,k_r}(\beta_1,\ldots,\beta_r)\notag\\
 &+\sum_{k=0}^{h-1}\sum_{j=0}^k
 b_{h,k}(-1)^j\binom kj\,
 \Li_{C-k+j,k_1-j,k_2,\ldots,k_r}
       (\alpha,\beta_1/\alpha,\beta_2,\ldots,\beta_r).
\end{align}
The first inner index $k_1-j$ is permitted to be nonpositive at this
intermediate step.
\end{lemma}
\begin{proof}
Convolution of coefficients gives a sum over $n\ge m\ge1$ with
factor $\alpha^{n-m}\binom{n-m+h-1}{h-1}$. The diagonal $n=m$
is the first term because the binomial coefficient equals one there.
For $n>m$, expand \eqref{mixed:b} and
$(n-m)^k=\sum_{j=0}^k(-1)^j\binom kj n^{k-j}m^j$.
The new outer color is $\alpha$ and the next is $\beta_1/\alpha$,
which proves the formula by direct reindexing.
\end{proof}

If $P=0$, no $H$ is present. The same binomial coefficient gives
directly
\begin{equation}\label{mixed:rationalonly}
 Z(C)=\sum_{\alpha,h}\sum_{k=0}^{h-1}
                    c_{\alpha,h}b_{h,k}\Li_{C-k}(\alpha),
\end{equation}
the depth-one theorem already proved in the source. Thus the new
construction agrees exactly with that existing sector.

\subsection{Step 3: eliminate every nonpositive inner slot}

\begin{lemma}[Finite elimination, including confluent colors]
\label{mixed:eliminate}
Suppose an inner index in \eqref{mixed:Li}, at position $j\ge2$,
equals $-M$ with $M\ge0$. Put $U=n_{j-1}$ and $L=n_{j+1}$;
for the final position set $L=0$. If $q=z_j\ne1$, define
\[
 R_h(q)=(q\,d/dq)^h\frac1{1-q}.
\]
Then
\begin{equation}\label{mixed:geometric}
 \sum_{n=L+1}^{U-1}q^n n^M
 =\sum_{h=0}^M\binom Mh R_h(q)
       \left\{q^{L+1}(L+1)^{M-h}-q^U U^{M-h}\right\}.
\end{equation}
At $q=1$ the exact replacement is
\begin{equation}\label{mixed:Bernoulli}
 \sum_{n=L+1}^{U-1}n^M
       =\frac{B_{M+1}(U)-B_{M+1}(L+1)}{M+1}.
\end{equation}
In either case the resulting multiple polylogarithms have one less
index. The cumulative alphabet does not enlarge.
\end{lemma}
\begin{proof}
Apply $(q\,d/dq)^M$ to the finite geometric sum
$(q^{L+1}-q^U)/(1-q)$ and use the product rule. This proves
\eqref{mixed:geometric}; the Bernoulli difference identity proves
\eqref{mixed:Bernoulli}.

Expand the powers of $L+1$ into powers of $L$. An upper-end term
multiplies the preceding color by $q$ and subtracts a nonnegative
integer from its order; a lower-end term does the same to the following
color and order. At the last position the lower-end term is a constant.
In every case the $j$th summation disappears. In cumulative colors,
the upper-end operation deletes the old $(j-1)$st cumulative letter,
and the lower-end operation deletes the old $j$th letter. When $q=1$
those two letters coincide. Hence no new cumulative color is introduced.
\end{proof}

If a neighboring positive index becomes nonpositive, repeat the rule.
Depth decreases at every repetition, so the algorithm terminates.
The first index can become any $C+$integer and is never eliminated.
This explains both the precise output class and the depth bound.

\begin{proof}[Completion of Theorem~\ref{mixed:main}]
Apply \eqref{mixed:roots}, the finite shuffle expansion, partial
fractions, Lemma~\ref{mixed:prefactor}, and
Lemma~\ref{mixed:eliminate}. The positive-only case has depth $d$;
a rational prefactor initially adds at most one index to depth $P$;
all subsequent eliminations decrease depth. The rational-only case is
\eqref{mixed:rationalonly}. Alphabet preservation and the root-order
statement follow at each stage. The proof so far uses locally normally
convergent series for interior cumulative colors, and absolute
convergence at unit colors in a sufficiently far right half-plane.

For analytic continuation use
\begin{equation}\label{mixed:Mellin}
 Z(C)=\frac1{\Gamma(C)}\int_0^\infty
                   t^{C-1}\prod_j\Li_{a_j}(X_je^{-p_jt})\dd t.
\end{equation}
It follows by termwise Laplace integration in a common convergent
half-plane. If no $X_j=1$, its kernel is analytic at zero and
exponentially decreasing at infinity. Subtracting an arbitrarily long
Taylor polynomial near zero gives meromorphic terms $1/(C+n)$,
whose poles are canceled by $1/\Gamma(C)$. This proves entire
continuation. The same argument applies to every output term: choose
an integral positive first index as its base Mellin kernel, and note
that its iterated-integral word has no singular letter at $t=1$.
Every cumulative letter lies in $\mathcal A$, which contains no $1$
when no original $X_j=1$. Its first-order continuation is therefore
entire. If some original $X_j=1$, the convergent power--logarithm
expansions below continue the product kernel meromorphically. Each
nested output term has the same type of continuation: a positive-index
iterated word has a convergent local power--logarithm expansion at
$t=1$. This follows recursively from its defining integrals, since
the only singular endpoint form is $\dd t/(1-t)$ and integrating a
convergent power series times a finite power of $\log(1-t)$ preserves
that class. Apply the same Mellin subtraction to a positive first-index
base word to continue its first order. The identity theorem extends
the finite reduction, and differentiation in the
free variable $C$ is legitimate.
\end{proof}

Radial continuation in the coefficient variable means $K(t)\mapsto
K(\rho t)$; in the original variables this is
$X_j\mapsto\rho^{p_j}X_j$, not a common damping of all original
colors. At resonant colors the resulting Abel limit need not be an
ordinary finite value outside the convergence domain. All formulas
there mean the stated meromorphic continuation, with Laurent
coefficients taken afterward.

\subsection{Local Laurent coefficients at every nonpositive outer order}

Let $P_1=\#\{j:a_j>0,\ X_j=1\}$. The integer limits of
Jonqui\`ere's expansion \cite[Eq. 25.12.12]{DLMFPolylog} give
\begin{align}\label{mixed:positiveexpansion}
 \Li_k(e^{-pt})
 ={}&\frac{(-pt)^{k-1}}{(k-1)!}
          \bigl(H_{k-1}-\log p-\log t\bigr)\notag\\
 &+\sum_{n\ge0,\ n\ne k-1}\zeta(k-n)\frac{(-pt)^n}{n!},
       \qquad k\ge1,\\
 \Li_{-M}(e^{-pt})
 ={}&\frac{M!}{(pt)^{M+1}}
          +\sum_{n\ge0}\zeta(-M-n)\frac{(-pt)^n}{n!},
       \qquad M\ge0.\label{mixed:negativeexpansion}
\end{align}
Here $H_0=0$. For $X\ne1$ the ordinary Taylor expansion is
\begin{equation}\label{mixed:colorexpansion}
 \Li_a(Xe^{-pt})=
       \sum_{n\ge0}\frac{(-p)^n}{n!}\Li_{a-n}(X)t^n.
\end{equation}
Products yield a convergent local power--logarithm expansion
\begin{equation}\label{mixed:local}
 \prod_j\Li_{a_j}(X_je^{-p_jt})
       =\sum_{n\ge n_0}\sum_{r=0}^{P_1}
                                  \kappa_{n,r}t^n(\log t)^r.
\end{equation}
At any prescribed $n$, its coefficients require only finitely many
terms of \eqref{mixed:positiveexpansion}--\eqref{mixed:colorexpansion}.

\begin{theorem}[Pole structure and finite local Laurent constants]
\label{mixed:localtheorem}
Set $g(C)=1/\Gamma(C)$. Near $C=-n$, the complete polar contribution
to $Z(C)$ is
\begin{equation}\label{mixed:localmodel}
        g(C)\sum_{r=0}^{P_1}
              \frac{(-1)^r r!\kappa_{n,r}}{(C+n)^{r+1}}.
\end{equation}
The difference from \eqref{mixed:localmodel} is $g(C)$ times a
holomorphic function. For $n\ge0$, in particular,
\begin{equation}\label{mixed:finitepart}
 \boxed{\displaystyle
 \CT_{C=-n} Z(C)=
     \sum_{r=0}^{P_1}\frac{(-1)^r r!}{(r+1)!}
                           \kappa_{n,r}g^{(r+1)}(-n).}
\end{equation}
Possible poles at nonpositive integers have order at most $P_1$;
possible poles at positive integers have order at most $P_1+1$.
There are no other possible poles. Missing local coefficients mean zero.
\end{theorem}
\begin{proof}
Taylor--logarithm subtraction in \eqref{mixed:Mellin} reduces the
local terms to
\[
 \int_0^1t^{C+n-1}(\log t)^r\dd t
                  =\frac{(-1)^r r!}{(C+n)^{r+1}}.
\]
Arbitrarily high subtraction makes the remainder holomorphic on any
prescribed compact set. At $C=-n$ with $n\ge0$, $g$ has a simple
zero. Thus its product with the holomorphic remainder has zero
constant coefficient, and the constant of each displayed local term
is exactly \eqref{mixed:finitepart}. The pole assertions follow from
the same zero and from the fact that the powers $n$ are integers.
\end{proof}

Formula \eqref{mixed:finitepart} does not determine all higher
regular coefficients. For example, the constant coefficient of
$Z'(C)$ generally receives the first coefficient of $g(C)$ times
the global holomorphic remainder. This is the same obstruction that
prevents differentiating an isolated negative-integer value identity
transversely. The finite mixed-order formula, which keeps $C$ free,
does determine such derivatives in its stated multiple-polylogarithm
class.

Reciprocal-Gamma derivatives in \eqref{mixed:finitepart} are finite
polynomials in $\gamma$, zeta values, and harmonic numbers. Indeed
\begin{equation}\label{mixed:reciprocalgamma}
 \frac1{\Gamma(-n+u)}=
       (-1)^n n!\,u\prod_{j=1}^n(1-u/j)
       \exp\left(\gamma u+
             \sum_{k\ge2}\frac{(-1)^{k+1}\zeta(k)}k u^k\right).
\end{equation}
Only degree $r+1$ is needed for a local logarithm of degree $r$.

\subsection{Explicit triple identities and a resonant Laurent expansion}

Define
\begin{equation}\label{mixed:F}
 F(C;z)=\sum_{m,n,k\ge1}\frac{m\,z^{m+n+k}}{nk(m+n+k)^C}.
\end{equation}
For $|z|<1$ it is entire in $C$; at $z=1$ the following original
series converges absolutely for $\Re C>2$. Its finite reduction is
\begin{align}\label{mixed:Fidentity}
 F(C;z)={}&2\Li_{C-1,1,1}(z,1,1)
          -2\Li_{C-1,1}(z,1)+2\Li_{C,1}(z,1)\notag\\
         &+2\Li_{C-1}(z)-2\Li_C(z).
\end{align}
An independent elementary proof groups by $N=m+n+k$:
\[
 \sum_{m+n+k=N}\frac m{nk}
 =2N H_{N-1}(1,1)-2(N-1)H_{N-1}+2(N-1).
\]
Here $H_N(1,1)=\sum_{N\ge j>k\ge1}(jk)^{-1}$.
The formula follows by first putting $v=n+k$, summing
$1/(n(v-n))=2H_{v-1}/v$, and using
$\sum_{v=1}^{N-1}H_{v-1}=(N-1)(H_{N-1}-1)$.
Thus this example tests the general construction independently of
the shuffle implementation.

The coefficient generator is
$\Li_{-1}(z)\Li_1(z)^2$. Consequently
\begin{align}\label{mixed:Fvalues}
 F(0;-1)&=-\frac{\log^2 2}{4},&
 F(-1;-1)&=-\frac{\log2}{4},\notag\\
 F(-2;-1)&=\frac{\log^2 2-\log2-1}{8},&
 F(1;-1)&=\frac{\log^2 2}{2}+\log2-1.
\end{align}
The first three follow by applying $(z\,d/dz)^{-C}$ to the rational
times logarithm generator at the nonpositive integers. For the last,
direct integration gives the stronger functional identity
\begin{equation}\label{mixed:Fprimitive}
 F(1;z)=\int_0^z\frac{\Li_1(u)^2}{(1-u)^2}\dd u
       =\frac{\Li_1(z)^2-2\Li_1(z)+2}{1-z}-2.
\end{equation}
More generally $z\,\partial_z F(C;z)=F(C-1;z)$ and the primitive
of $F(C;z)/z$ vanishing at zero is $F(C+1;z)$. Every fixed
$C$-derivative obeys the same identities.

At $z=1$, the kernel in \eqref{mixed:Mellin} is
\[
 \Li_{-1}(e^{-t})\Li_1(e^{-t})^2
 =\frac{\log^2t}{t^2}-\frac{\log t}{t}
    -\frac{\log^2t}{12}+\frac{\log t}{12}+\frac14
       +O\bigl(t(1+|\log t|^2)\bigr).
\]
Using \eqref{mixed:localmodel} and \eqref{mixed:reciprocalgamma}
therefore proves the exact Laurent expansion
\begin{equation}\label{mixed:Flaurent}
 \boxed{\displaystyle
 F(C;1)=-\frac1{6C^2}-\frac{2\gamma+1}{12C}
       +\frac14-\frac{\gamma^2+\gamma}{12}
                +\frac{\zeta(2)}{12}+O(C).}
\end{equation}
This example has a genuine double pole at a nonpositive outer order.
The finite part is a short Gamma--zeta polynomial, while the whole
function is the depth-three expression \eqref{mixed:Fidentity}.

A second example exercises two different nonpositive-slot eliminations:
\begin{align}\label{mixed:squareexample}
 \sum_{m,n,k\ge1}\frac{k^2z^{m+n+k}}{mn(m+n+k)^C}
 ={}&2\Li_{C-2,1,1}(z,1,1)
       -3\Li_{C-2,1}(z,1)+3\Li_{C-1,1}(z,1)\notag\\
 &+\frac72\Li_{C-2}(z)-\frac92\Li_{C-1}(z)+\Li_C(z).
\end{align}
Its generator is $\Li_{-2}(z)\Li_1(z)^2$. At $z=1$ its finite
value at $C=0$ is $-1/12$. At $C=-1+u$ its Laurent expansion is
\begin{equation}\label{mixed:squareLaurent}
 \frac1{60u^2}+\frac{\gamma/60-13/720}{u}
   +\frac{\gamma^2}{120}-\frac{13\gamma}{720}
    -\frac{\pi^2}{720}-\frac1{480}+O(u).
\end{equation}
For verification, the $t^1$ coefficient of the local kernel is
$-\log^2t/120-\log t/720+1/288$; insertion into
\eqref{mixed:localmodel} at $C=-1$ gives
\eqref{mixed:squareLaurent} directly.

\subsection{A convergent evaluation and exact pole classification}

Write $U(C;z)$ for the left side of \eqref{mixed:squareexample}.
At $z=1$ its defining series converges absolutely for $\Re C>3$.
Grouping by $N=m+n+k$ gives the additional finite harmonic identity
\begin{align}\label{mixed:Ucoeff}
 \sum_{m+n+k=N}\frac{k^2}{mn}
 ={}&N^2\bigl(H_{N-1}^2-H_{N-1}^{(2)}\bigr)
        -3N(N-1)H_{N-1}\notag\\
 &+\frac{7N^2-9N+2}{2}.
\end{align}
This follows by taking the outer coefficient in
\eqref{mixed:squareexample}, including the vanishing cases $N=1,2$.
The following ordinary value uses only classical height-one multiple
zeta identities, whose beta-integral proof is included to fix the
normalization.

\begin{corollary}[A convergent zeta evaluation]
\label{mxex:weight-four}
\begin{equation}
 \boxed{\displaystyle
 \sum_{m,n,k\geq1}\frac{k^2}{mn(m+n+k)^4}
   =\frac{\pi^4}{24}+\frac{7\pi^2}{12}
      -\frac{15}{2}\zeta(3).}
 \label{mxex:convergent-value}
\end{equation}
\end{corollary}

\begin{proof}
Set $C=4,z=1$ in \eqref{mixed:squareexample}, and use
\[
 \zeta(2,1,1)=\zeta(4),\qquad
 \zeta(2,1)=\zeta(3),\qquad
 \zeta(3,1)=\frac14\zeta(4).
\]
For a self-contained verification of these three classical relations,
consider the classical Aomoto--Drinfel'd height-one generating identity,
recorded in Kaneko and Ohno~\cite[Eq.~(5)]{mxex:KanekoOhno2010}:
\begin{equation}\label{mxex:height-one}
 \sum_{r,s\geq1}\zeta(r+1,\{1\}^{s-1})x^ry^s
   =1-\frac{\Gamma(1-x)\Gamma(1-y)}{\Gamma(1-x-y)}.
\end{equation}
Indeed the shuffle identity gives
$(-\log(1-t))^s/s!=\operatorname{Li}_{1,\ldots,1}(t,1,\ldots,1)$,
with $s$ indices.  Expand this absolutely convergent power series and
use
$\int_0^1t^{n-1}(-\log t)^{r-1}\,dt=(r-1)!/n^r$.
Termwise integration is justified by nonnegative summands and gives
\[
 \zeta(r+1,\{1\}^{s-1})
 =\frac1{(r-1)!s!}\int_0^1
   \frac{(-\log t)^{r-1}(-\log(1-t))^s}{t}\,dt.
\]
Summing on a sufficiently small bidisc turns the right side into
\[
 x\int_0^1t^{-x-1}\bigl((1-t)^{-y}-1\bigr)\,dt
 =1+x\frac{\Gamma(-x)\Gamma(1-y)}{\Gamma(1-x-y)}.
\]
The beta evaluation is first justified for $\Re x<0$ and then by
analytic continuation near zero.  The functional equation of Gamma
gives \eqref{mxex:height-one}.  Finally
\[
 \log\frac{\Gamma(1-x)\Gamma(1-y)}{\Gamma(1-x-y)}
 =\sum_{j\geq2}\frac{\zeta(j)}j
   \bigl(x^j+y^j-(x+y)^j\bigr).
\]
Its $xy^2$, $xy^3$, and $x^2y^2$ coefficients give the three displayed
identities; in the last case use $\zeta(2)^2=5\zeta(4)/2$.
Formula \eqref{mixed:squareexample} now gives
$15\zeta(4)/4+7\zeta(2)/2-15\zeta(3)/2$, which is
\eqref{mxex:convergent-value}.
\end{proof}

\subsection{The full pole pattern at the resonant color}

Write $U(C)=U(C;1)$ and let $B_j$ denote the Bernoulli numbers with
$B_1=-1/2$.  For $\Re C>3$,
\begin{equation}\label{mxex:mellin}
 U(C)=\frac1{\Gamma(C)}\int_0^\infty t^{C-1}K(t)\,dt,
 \qquad
 K(t)=\operatorname{Li}_1(e^{-t})^2
                  \operatorname{Li}_{-2}(e^{-t}).
\end{equation}
The integral follows directly from the Gamma integral for the denominator;
absolute convergence follows from the coefficient bound $O(N^2(1+\log N)^2)$.

\begin{proposition}[Actual spectral poles]
\label{mxex:poles}
The meromorphic function $U$ has a triple pole at $C=3$, double poles
at $C=2,1$ and at every negative odd integer, and simple poles at
every negative even integer.  These are all its poles.  The point
$C=0$ is regular and
\begin{equation}\label{mxex:zero-value}
 U(0)=-\frac1{12}.
\end{equation}
For every integer $m\geq1$,
\begin{equation}\label{mxex:even-residue}
 \operatorname*{Res}_{C=-2m}U(C)
       =\frac{m}{m+1}B_{2m+2}.
\end{equation}
For every integer $m\geq0$, the coefficient of the leading double
pole at $C=-(2m+1)$ is
\begin{equation}\label{mxex:odd-leading}
 [(C+2m+1)^{-2}]\,U(C)
       =-\frac{B_{2m+4}}{m+2}.
\end{equation}
\end{proposition}

\begin{proof}
The convergent local expansions for $0<t<2\pi$ are
\begin{align*}
 \operatorname{Li}_1(e^{-t})&=-\log t+a(t),&
 a(t)&=\frac t2-\sum_{r\geq1}
             \frac{B_{2r}}{2r(2r)!}t^{2r},\\
 \operatorname{Li}_{-2}(e^{-t})&=b(t),&
 b(t)&=\frac2{t^3}+\sum_{r\geq1}
             \frac{B_{2r+2}}{(2r+2)(2r-1)!}t^{2r-1}.
\end{align*}
The second expansion follows by differentiating
$(e^t-1)^{-1}$ twice.  The first follows by integrating its negative
and fixing the constant from
$\operatorname{Li}_1(e^{-t})+\log t\to0$.
Thus
\[
 K(t)=b(t)(\log t)^2-2a(t)b(t)\log t+a(t)^2b(t).
\]
Mellin subtraction shows that a coefficient
$\kappa_{n,r}t^n(\log t)^r$ contributes the singular model
\[
 \frac1{\Gamma(C)}
 \frac{(-1)^r r!\kappa_{n,r}}{(C+n)^{r+1}}.
\]
All other contributions are holomorphic locally, after a sufficiently
long Taylor--logarithm subtraction.  There are no additional poles
because the tail of the integral decays exponentially.

At the first four powers,
\begin{equation}\label{mxex:local-first}
 K(t)=\frac{2\log^2t}{t^3}
       -\frac{2\log t}{t^2}
       +\frac{\log t/6+1/2}{t}
       -\frac1{12}+O(t\log^2t).
\end{equation}
The three positive poles therefore have orders three, two, and two,
with nonzero leading coefficients $2$, $2$, and $-1/6$, respectively.
The zero of $1/\Gamma(C)$ at $C=0$ cancels the simple local Mellin
pole there; its derivative at zero is one, proving
\eqref{mxex:zero-value}.

For a positive even $n$, $b(t)$ has no $t^n$ term.  Since the only
odd power in $a(t)$ is $t/2$, the coefficient of $t^n\log t$ in
$K$ is $-b_{n-1}$, where
\[
 b_{n-1}=\frac{B_{n+2}}{(n+2)(n-1)!}.
\]
Using $(1/\Gamma)'(-n)=(-1)^n n!$ gives residue
$n!b_{n-1}=nB_{n+2}/(n+2)$, which is
\eqref{mxex:even-residue}.  For positive odd $n$, the coefficient
of $t^n\log^2t$ is
$b_n=B_{n+3}/((n+3)n!)$.  The same reciprocal-Gamma zero gives
leading double-pole coefficient $-2n!b_n$, proving
\eqref{mxex:odd-leading}.  The even-index Bernoulli numbers appearing
here are nonzero, so these are actual poles, of exactly the stated
orders.
\end{proof}


At the next nonpositive even order one also obtains
\begin{equation}\label{mixed:Uminus2}
 \CT_{C=-2}U(C)=\frac{19}{720}-\frac{\gamma}{60}.
\end{equation}
Indeed $[t^2]K(t)=\log(t)/120+1/1440$, and
$1/\Gamma(-2+u)=2u+(2\gamma-3)u^2+O(u^3)$.
Multiplication by its local Mellin model gives
$2/1440-(2\gamma-3)/120$, proving the displayed constant.
This is a Laurent constant of the continuation, not a value of the
positive-term defining series at $C=-2$.

\subsection{The depth-one boundary: Gamma, polygamma, and Stieltjes data}

When the normal form reaches depth one at a $q$th root of unity
$\alpha$, finite residue-class summation gives
\begin{equation}\label{mixed:DFT}
 \Li_s(\alpha)=q^{-s}\sum_{j=1}^q\alpha^j\zeta(s,j/q).
\end{equation}
For $s\ne1$ its $k$th order derivative is
\begin{equation}\label{mixed:DFTjets}
 q^{-s}\sum_{j=1}^q\alpha^j
  \sum_{\ell=0}^k\binom k\ell(-\log q)^{k-\ell}
                         \zeta^{(\ell)}(s,j/q).
\end{equation}
For nontrivial $\alpha$ the common pole at $s=1$ cancels first,
so the finite formula there is
\begin{equation}\label{mixed:DFTStieltjes}
 \Li_1^{[k]}(\alpha)=\frac1q\sum_{j=1}^q\alpha^j
  \sum_{\ell=0}^k\binom k\ell(-\log q)^{k-\ell}
                         (-1)^\ell\gamma_\ell(j/q).
\end{equation}
These classical identities, already used in the source, fix the
normalization of the depth-one output. Lerch's identity
$\zeta'(0,a)=\log\Gamma(a)-\tfrac12\log(2\pi)$ and
$\psi^{(r)}(a)=(-1)^{r+1}r!\zeta(r+1,a)$, $r\ge1$, give the corresponding
Gamma and polygamma expressions \cite{DLMFHurwitz}. The new theorem
does not erase the positive-depth order derivatives left by its
normal form; further reductions require additional identities.
